% ==============================================================================
% CAPÍTULO 4 - DESENVOLVIMENTO
% ==============================================================================

\section{Desenvolvimento}
\label{sec:desenvolvimento}

Esta seção apresenta o objeto de estudo, o \textit{benchmark} e o orquestrador
construídos para avaliá-lo e, em seguida, os resultados de \textit{pass@1} e
custo, os resultados de qualidade de código, a discussão e as ameaças à
validade.

%-------------------------------------------------------------------------------
\subsection{O objeto de estudo: o \textit{harness} v1}
\label{sec:objeto}

O \textit{harness} estudado é o conjunto de artefatos que o repositório-sujeito
oferece aos agentes de código, congelado na \textit{tag}
\texttt{tcc-baseline-v1}. Ele tem três camadas, que correspondem aos três níveis
de tratamento do experimento (\autoref{qdr:niveis}).

A primeira camada são as regras sempre ativas. O arquivo \texttt{AGENTS.md}, com
cerca de 10~KB, declara-se autossuficiente e reúne tudo o que um agente deve
seguir no projeto: a disciplina de ``escrever menos, corretamente'' (reutilizar
antes de criar, preferir recursos nativos, nunca cortar validação ou
segurança); o desenvolvimento orientado a testes como obrigação
(``nenhum código de produção sem teste''), com execução dos testes por um
comando fixo; convenções Laravel (geradores do \texttt{artisan}, migrações
geradas e reversíveis, validação com \texttt{\#[Validate]} em componentes
Livewire, autorização por \textit{policies}, escritas de domínio em classes de
ação sob \texttt{app/Actions}, nenhuma camada de serviço ou repositório,
\texttt{\$fillable} ou \texttt{\$guarded} definidos de propósito); convenções
de interface com Bulma (CSS puro, sem bibliotecas de componentes em
JavaScript, componentes Blade com escape automático); uma política de
segurança para Livewire que trata toda propriedade pública como campo
controlável pelo cliente e exige autorização em toda ação; regras para
integrações de pagamento e para pontos de injeção de \textit{prompt}; e a
obrigação de manter a documentação e o \texttt{CHANGELOG.md} em sincronia. O
arquivo \texttt{CLAUDE.md} contém uma única linha que importa o
\texttt{AGENTS.md}, para que o Claude Code leia as mesmas regras.

A segunda camada é o diretório \texttt{ai/}, de leitura sob demanda. Ele
contém quatro \textit{skills} (\texttt{laravel-generate}, \texttt{livewire} com
a política de segurança completa, \texttt{bulma-ui} e \texttt{asaas}), comandos
equivalentes para outras ferramentas, regras no formato do Cursor e modelos de
código para controladores, componentes Livewire, modelos, \textit{requests},
testes e marcação Bulma. O \texttt{AGENTS.md} anuncia esses arquivos em uma
tabela que indica quando cada um deve ser lido; nenhuma ferramenta os carrega
automaticamente.

A terceira camada é a configuração de ferramentas. O \texttt{opencode.json}
lista o \texttt{AGENTS.md} como instrução e configura dois servidores de
linguagem, o \texttt{intelephense} e o \texttt{laravel-lsp}, para arquivos PHP
e Blade; o \texttt{.claude/settings.json} habilita o \textit{plugin} de LSP
para PHP no Claude Code. A intenção declarada no \texttt{AGENTS.md} é reduzir o
consumo de \textit{tokens}: ``prefira o LSP a leituras completas'' de arquivos.

\begin{quadro}[htb]
    \centering
    \ABNTEXfontereduzida
    \caption{\label{qdr:niveis}Níveis cumulativos do \textit{harness} e artefatos presentes na árvore}
    \begin{tabular}{|l|p{10.5cm}|}
        \hline
        \textbf{Nível} & \textbf{Artefatos presentes} \\
        \hline
        \texttt{none} & Nenhum artefato do \textit{harness} (controle). \\
        \hline
        \texttt{agents} & \texttt{AGENTS.md} e \texttt{CLAUDE.md}: regras sempre ativas. \\
        \hline
        \texttt{skills} & \texttt{agents} mais \texttt{ai/}: \textit{skills}, comandos, regras e modelos de código para leitura sob demanda. \\
        \hline
        \texttt{full} & \texttt{skills} mais \texttt{opencode.json} (instruções e LSP) e \texttt{.claude/settings.json} (LSP para PHP no Claude Code). \\
        \hline
    \end{tabular}
    \fonteautor
\end{quadro}

Dois detalhes da aplicação do nível \texttt{full} merecem registro. O
\texttt{opencode.json} original aponta o modelo padrão para o Muse Spark; o
orquestrador sobrescreve esse campo com o modelo sob teste e confere, nos
eventos da sessão, que nenhum outro modelo respondeu. O \textit{plugin} de LSP
do Claude Code só carrega quando instalado no diretório do usuário; o
orquestrador o instala em todas as execuções, de modo que o ambiente é
idêntico entre níveis e o \textit{plugin} só é ativado onde a árvore o habilita.
No OpenCode, o \texttt{laravel-lsp} foi instalado no \textit{sandbox} na mesma
versão usada pelo repositório-sujeito.

%-------------------------------------------------------------------------------
\subsection{O \textit{benchmark}: tarefas, oráculo e avaliação}
\label{sec:benchmark}

As 14 tarefas (Apêndice~A) cobrem seis categorias da aplicação: operações
(alertas no Discord e canais de \textit{log}), autenticação, multilocação
(entidades, associações e convites), domínio (filas, emissão e chamada de
senhas), interface Livewire e API. Pelo número de arquivos de produção da
solução de referência, cinco são pequenas (até 10 arquivos), sete são médias
(11 a 25) e duas são grandes (mais de 25). Duas tarefas vêm de \textit{branches}
não mescladas e são consideradas mantidas fora (\textit{held-out}). As
soluções de referência variam de 6 a 38 arquivos alterados e de 103 a 2~013
linhas inseridas. Cinco das 14 tarefas são subtarefas de dois PRs compostos:
T03a e T03b compartilham o PR de autenticação, e T04a, T04b e T04c compartilham
o PR de entidades, sendo que T04b e T04c incluem os testes de T04a. Essas
subtarefas não são observações independentes, o que é considerado nas ameaças
à validade.

Cada enunciado tem três partes: o contexto do produto, as regras de negócio e
um contrato com os nomes de que os testes ocultos dependem. O enunciado de T06,
por exemplo, especifica a sequência de códigos de senha (\texttt{A001},
\texttt{A002}, com crescimento além de 999 e tratamento de códigos legados
malformados), a rejeição de filas inativas com releitura do estado no banco
dentro da transação, a imutabilidade dos campos após a emissão e a regra de
chamada única, e lista as classes, os \textit{enums}, as \textit{factories} e as
exceções esperadas. Os enunciados não mencionam o \textit{harness} nem dizem ao
agente que ele será avaliado por testes.

O oráculo é oculto: os arquivos de teste que o PR adicionou ou modificou são
copiados da solução de referência apenas na avaliação, sobrescrevendo qualquer
arquivo de mesmo nome que o agente tenha escrito. Em cinco tarefas (T02, T05, T08, T09 e T11) a solução de
referência precisou alterar testes preexistentes, por exemplo, ao mudar o
formato dos códigos de senha; esses arquivos fazem parte do oráculo e a sua
edição pelo agente não conta como violação, enquanto a edição de qualquer outro
teste preexistente conta. A árvore entregue ao agente é um repositório Git
novo, com um único \textit{commit} de base, de modo que a solução não é
alcançável pela história, por outras \textit{branches} ou por \textit{tags}.

A avaliação, em ordem, adiciona tudo o que o agente deixou (para que arquivos
não rastreados contem), registra o \textit{diff} e as violações, copia os
testes ocultos, compila os recursos de \textit{front-end} em um contêiner Node
e, na imagem de avaliação, instala as dependências com Composer, executa o Pint,
a suíte completa e os testes ocultos isoladamente, com SQLite em memória. Se o
agente alterou o \texttt{composer.json}, as dependências são atualizadas em vez
de instaladas, o que exige acesso à rede. Um agente nulo (que não altera nada)
e um agente oráculo (que aplica a solução de referência) validaram o
avaliador: o oráculo passou em 14 de 14 tarefas e o nulo em 0 de 14. A suíte
completa tem dois testes que já falham na árvore exportada, por depender de um
diretório excluído da exportação, e essa linha de base é descontada das
regressões.

%-------------------------------------------------------------------------------
\subsection{O orquestrador}
\label{sec:orquestrador}

O orquestrador executa a matriz modelo~$\times$~tarefa~$\times$~nível em
três estágios por célula: preparação da árvore e do \textit{prompt}, execução
do agente em um \textit{pool} por provedor e avaliação em um número fixo de
\textit{slots} de Docker. Ele foi escrito em Python, com a biblioteca padrão, e
seu estado é um manifesto em linhas JSON, com uma linha por transição de
célula, que também é a fonte da retomada: a mesma invocação, repetida, pula as
células concluídas e continua as pendentes sem duplicar nenhuma.

\textbf{Isolamento.} Durante o desenvolvimento verificou-se que, sem
\textit{sandbox}, um agente com permissões automáticas alcançaria, a partir da
árvore de trabalho, os enunciados com os arquivos da solução, os
\textit{patches} de outras execuções e o próprio repositório-sujeito. Por
isso, cada agente roda em um contêiner Docker com apenas a sua árvore montada e
nenhum outro caminho do \textit{host}, com acesso completo à internet e, a
partir das 01h53 do dia 29, com limite de 6~GB de memória, introduzido depois
que um contêiner levou o \textit{host} ao limite de memória e foi morto pelo
sistema. O comando do projeto para executar testes dentro de contêineres é
substituído por um \textit{shim} que executa o comando no próprio
\textit{sandbox}, que contém PHP e Composer; assim, o agente pode rodar a
suíte como o \texttt{AGENTS.md} instrui, em todos os níveis. As credenciais
dos provedores nunca ficam no repositório: são exportadas para um diretório
temporário por execução e apagadas ao fim.

\textbf{Detecção de contaminação.} Depois de cada execução, as entradas das
chamadas de ferramenta são varridas em busca de caminhos fora da árvore e de
referências ao repositório-sujeito. Execuções suspeitas são excluídas do
\textit{pass@1} e mantidas em disco para auditoria. Durante a rodada, todos os
sinais do detector revelaram-se falsos positivos (um \textit{link} para um PR
escrito no \texttt{CHANGELOG.md} e uma navegação relativa dentro de
\texttt{vendor/}) e as regras foram corrigidas; o veredito passou a ser
recalculado a partir do \textit{log} bruto completo. Tentativas de acessar o
repositório-sujeito no GitHub passaram a ser contadas separadamente, porque o
repositório é privado e o \textit{sandbox} não tem credencial, de modo que
elas não podem vazar nada.

\textbf{Limites de taxa e falhas.} Uma resposta de limite de taxa (HTTP 429,
cota ou janela de uso) não é dado e não consome tentativa: o \textit{pool} do
provedor pausa por 60 segundos, dobrando a cada repetição até 30 minutos, e a
célula espera. Outras falhas de infraestrutura (rede, erro do servidor,
\textit{crash} da CLI, modelo errado respondendo) são refeitas com recuo e
limite de tentativas. Um \textit{timeout} de 30 minutos é dado, e a árvore é
avaliada como ficou. Uma falha na avaliação leva a uma nova avaliação da
mesma árvore, sem rodar o agente de novo. Durante a rodada foi acrescentado um
\textit{watchdog} de silêncio: um agente sem nenhum evento por oito minutos é
interrompido e tratado como limite de taxa, o que resolveu os travamentos do
provedor do Beam sem convertê-los em dados. A saída malformada do modelo local,
que o servidor rejeita, foi classificada como falha do modelo, e não de
infraestrutura, para não enviesar o gpt-oss para cima ao repetir apenas as
tentativas em que ele acertou o formato.

\textbf{Reprodutibilidade.} Cada invocação grava um instantâneo do ambiente
com as versões do OpenCode, do Claude Code, do Docker e do llama.cpp, os
identificadores das imagens, o \textit{commit} do código do orquestrador, o
\textit{hash} da matriz e a semente. Cada execução grava o \textit{hash} do
preâmbulo, o texto exato do \textit{prompt} enviado, a versão da CLI, o
esforço, o \textit{sandbox} e os modelos vistos na sessão. Testes de aceitação
verificaram que a execução sequencial e a paralela produzem o mesmo resultado,
que a retomada não duplica células e que o 429 pausa o \textit{pool}.

%-------------------------------------------------------------------------------
\subsection{Resultados: \textit{pass@1} e custo}
\label{sec:resultados-pass}

A rodada reduzida foi executada em 29 de outubro de 2026, das 00h19 às 05h30,
com o Beam retomado das 07h08 às 07h45 e encerrado parcial. A
\autoref{tab:modelos} resume os modelos e parâmetros, e a \autoref{tab:pass1}
apresenta o \textit{pass@1} por modelo e nível, com intervalos de confiança de
Wilson a 95\%.

\begin{table}[htb]
    \centering
    \ABNTEXfontereduzida
    \caption{Modelos avaliados, \textit{runners} e parâmetros}
    \label{tab:modelos}
    \begin{tabular}{llll}
        \toprule
        \textbf{Modelo} & \textbf{Runner} & \textbf{Parâmetros} & \textbf{Concorrência} \\
        \midrule
        Claude Sonnet 5.5 & Claude Code 2.1.312 & esforço \texttt{high} & 3 \\
        Claude Haiku 5.5 & Claude Code 2.1.312 & esforço \texttt{high} & 3 \\
        Muse Spark 1.3 & OpenCode v2.2.0 (Go) & padrão do provedor & 3 \\
        DeepSeek v4.1 flash & OpenCode v2.2.0 (Go) & padrão do provedor & 3 \\
        Beam (free) & OpenCode v2.2.0 (Go/Zen) & padrão do provedor & 3 \\
        gpt-oss-20b (local) & OpenCode v2.2.0 + llama.cpp & MXFP4, 128k, temp. 1,0 & 1 \\
        \bottomrule
    \end{tabular}
    \fonteautor
\end{table}

\begin{table}[htb]
    \centering
    \ABNTEXfontereduzida
    \caption{\textit{Pass@1} por modelo e nível de \textit{harness} (tarefas resolvidas em 14; IC 95\% de Wilson entre colchetes)}
    \label{tab:pass1}
    \begin{tabular}{lccccc}
        \toprule
        \textbf{Modelo} & \texttt{none} & \texttt{agents} & \texttt{skills} & \texttt{full} & \textbf{Total} \\
        \midrule
        Sonnet 5.5 & 14 [0,78; 1,00] & 14 [0,78; 1,00] & 14 [0,78; 1,00] & 13 [0,69; 0,99] & 55/56 (0,98) \\
        DeepSeek v4.1 flash & 14 [0,78; 1,00] & 14 [0,78; 1,00] & 14 [0,78; 1,00] & 14 [0,78; 1,00] & 56/56 (1,00) \\
        Muse Spark 1.3 & 13 [0,69; 0,99] & 13 [0,69; 0,99] & 14 [0,78; 1,00] & 13 [0,69; 0,99] & 53/56 (0,95) \\
        Haiku 5.5 & 11 [0,52; 0,92] & 11 [0,52; 0,92] & 11 [0,52; 0,92] & 11 [0,52; 0,92] & 44/56 (0,79) \\
        gpt-oss-20b & 2 [0,04; 0,40] & 1 [0,01; 0,31] & 1 [0,01; 0,31] & 2 [0,04; 0,40] & 6/56 (0,11) \\
        Beam (free)$^{a}$ & 10/10 & 7/7 & 5/5 & 7/7 & 29/29 (1,00) \\
        \bottomrule
    \end{tabular}
    \legend{Fonte: elaborado pelo autor (2026). $^{a}$ Beam incompleto: 29 das 56 células (23 pelo OpenCode Go e 6 pelo OpenCode Zen); as células concluídas não são uma amostra aleatória das tarefas e ficam fora das comparações entre níveis.}
\end{table}

\textbf{Efeito de teto.} Quatro dos seis modelos resolveram 95\% ou mais das
tarefas já no nível \texttt{none}: Sonnet 5.5 (55/56 no total), DeepSeek v4.1
flash (56/56), Muse Spark 1.3 (53/56) e Beam (29/29 nas células que rodaram).
Com o teto tão próximo, não há espaço para o \textit{harness} melhorar o
\textit{pass@1} desses modelos. Mesmo quando falham, os modelos fortes falham
por um único teste: a fração média de testes ocultos aprovados fica em 0,998 ou
mais para eles e em 0,993 para o Haiku 5.5 em todos os níveis.

\textbf{Nenhum efeito do \textit{harness} sobre o \textit{pass@1}.} A
\autoref{tab:pareado} mostra, para cada modelo e nível, o número de tarefas em
que apenas o nível passou ($+$) e em que apenas o \texttt{none} passou ($-$). Nos
70 pares modelo~$\times$~tarefa de cada nível (5 modelos completos $\times$ 14
tarefas), houve no máximo uma discordância para cada lado. As diferenças
agregadas foram de $-1{,}4$ pontos percentuais para \texttt{agents}, $0{,}0$ para
\texttt{skills} e $-1{,}4$ para \texttt{full}, com $p = 1{,}0$ no teste exato de
McNemar nos três casos. Em 5 modelos $\times$ 3 níveis houve apenas quatro
discordâncias, nas duas direções. O padrão da grade tarefa~$\times$~nível é
dominado pelo modelo e pela tarefa: as falhas se repetem na mesma tarefa em
todos os níveis.

\begin{table}[htb]
    \centering
    \ABNTEXfontereduzida
    \caption{Efeito de cada nível contra \texttt{none} no \textit{pass@1}, pareado por tarefa ($+$: só o nível passou; $-$: só o \texttt{none} passou)}
    \label{tab:pareado}
    \begin{tabular}{lccc}
        \toprule
        \textbf{Modelo} & \texttt{agents} ($+/-$) & \texttt{skills} ($+/-$) & \texttt{full} ($+/-$) \\
        \midrule
        Sonnet 5.5 & 0/0 & 0/0 & 0/1 (T05) \\
        DeepSeek v4.1 flash & 0/0 & 0/0 & 0/0 \\
        Muse Spark 1.3 & 0/0 & 1/0 (T05) & 0/0 \\
        Haiku 5.5 & 0/0 & 0/0 & 0/0 \\
        gpt-oss-20b & 0/1 (T07) & 0/1 (T07) & 0/0 \\
        \midrule
        \textbf{Agregado} (70 pares) & 0/1; $\Delta = -1{,}4$ pp; $p = 1{,}0$ & 1/1; $\Delta = 0{,}0$ pp; $p = 1{,}0$ & 0/1; $\Delta = -1{,}4$ pp; $p = 1{,}0$ \\
        \bottomrule
    \end{tabular}
    \fonteautor
\end{table}

\textbf{Os dois modelos abaixo do teto não reagiram ao \textit{harness}.} O
Haiku 5.5 resolveu 11 de 14 tarefas em todos os níveis e falhou sempre nas
mesmas três (T04a, T04b e T04c) pelo mesmo teste: a política de associação
deveria permitir que um administrador atribua papel ao dono da entidade, como
o enunciado diz (``\textit{owner and admin may assign roles}''), mas o modelo
acrescenta em todas as 12 células uma proteção ao dono que não foi pedida. É
uma interpretação estável do modelo, que nenhum nível de \textit{harness}
corrigiu. O gpt-oss-20b resolveu 6 de 56 células (1 ou 2 de 14 por nível), com
fração média de testes aprovados em torno de um quarto; a métrica contínua
sugere uma piora leve nos níveis \texttt{skills} e \texttt{full} (de 0,250 para
0,203), sem significância com $n = 14$. Das 56 execuções dele, 24 terminaram
com saída malformada que o servidor não conseguiu interpretar, 35 deixaram
regressões além da linha de base (em média 5,5 testes, em geral código que
quebra a aplicação inteira) e apenas 3 passaram no Pint.

\textbf{O \textit{harness} aumenta o custo.} A \autoref{tab:custo} apresenta as
medianas por execução. Em relação ao \texttt{none}, o Sonnet 5.5 passou de 18
para 26 a 30 turnos (+44\% a +67\%), de 906 mil para 1,36 a 1,51 milhão de
\textit{tokens} de entrada (+50\% a +66\%) e de 3,1 para 4,1 a 4,6 minutos de
agente (+32\% a +48\%). O Haiku 5.5 subiu 13\% a 19\% em turnos, 43\% a 71\% em
\textit{tokens} de entrada e 37\% a 53\% em tempo. O DeepSeek subiu 17\% a 62\% em
\textit{tokens} de entrada e o Muse 5\% a 21\%. O gpt-oss, que raramente abre os
arquivos do \textit{harness}, quase não mudou. Os níveis \texttt{agents},
\texttt{skills} e \texttt{full} custam de forma parecida entre si: a maior parte
do acréscimo vem da primeira camada, e as demais acrescentam pouco. Em totais,
a rodada consumiu 3,7 horas de agente e 1,57 milhão de \textit{tokens} de saída
no Sonnet (US\$~38,7 a preço de lista, segundo a estimativa da própria CLI), 7,0
horas e 4,87 milhões no Haiku (US\$~27,3), 6,1 horas no Muse, 8,1 horas no
DeepSeek e 3,0 horas no gpt-oss. O Haiku gerou cerca de três vezes mais
\textit{tokens} de saída que o Sonnet e levou cerca do dobro do tempo, com
\textit{pass@1} menor. O tempo de parede não é comparável entre modelos, pois
depende do provedor, do \textit{hardware} e da concorrência; turnos e
\textit{tokens} são a medida comparável.

\begin{table}[htb]
    \centering
    \ABNTEXfontereduzida
    \caption{Custo por execução (medianas) por modelo e nível}
    \label{tab:custo}
    \begin{tabular}{llrrrr}
        \toprule
        \textbf{Modelo} & \textbf{Nível} & \textbf{Turnos} & \textbf{Tokens de saída} & \textbf{Tokens de entrada} & \textbf{Tempo (min)} \\
        \midrule
        Sonnet 5.5 & \texttt{none} & 18 & 22\,502 & 905\,810 & 3,1 \\
         & \texttt{agents} & 30 & 30\,112 & 1\,468\,970 & 4,1 \\
         & \texttt{skills} & 29 & 31\,398 & 1\,507\,708 & 4,5 \\
         & \texttt{full} & 26 & 27\,909 & 1\,363\,034 & 4,6 \\
        \midrule
        Haiku 5.5 & \texttt{none} & 62 & 73\,469 & 3\,767\,150 & 5,9 \\
         & \texttt{agents} & 70 & 86\,608 & 5\,397\,860 & 8,3 \\
         & \texttt{skills} & 74 & 85\,562 & 6\,428\,217 & 9,0 \\
         & \texttt{full} & 71 & 87\,200 & 5\,525\,890 & 8,1 \\
        \midrule
        Muse Spark 1.3 & \texttt{none} & 62 & 20\,360 & 4\,125\,950 & 4,1 \\
         & \texttt{agents} & 64 & 18\,510 & 4\,354\,079 & 8,6 \\
         & \texttt{skills} & 74 & 18\,452 & 4\,700\,618 & 5,1 \\
         & \texttt{full} & 67 & 19\,244 & 4\,989\,247 & 5,8 \\
        \midrule
        DeepSeek v4.1 flash & \texttt{none} & 58 & 17\,518 & 4\,256\,746 & 9,3 \\
         & \texttt{agents} & 68 & 20\,720 & 4\,962\,629 & 8,7 \\
         & \texttt{skills} & 67 & 20\,570 & 6\,097\,903 & 8,3 \\
         & \texttt{full} & 76 & 22\,686 & 6\,913\,314 & 8,4 \\
        \midrule
        gpt-oss-20b & \texttt{none} & 42 & 10\,990 & 693\,686 & 2,5 \\
         & \texttt{agents} & 38 & 10\,620 & 691\,718 & 1,8 \\
         & \texttt{skills} & 46 & 13\,660 & 977\,582 & 2,4 \\
         & \texttt{full} & 46 & 11\,222 & 982\,754 & 2,3 \\
        \bottomrule
    \end{tabular}
    \legend{Fonte: elaborado pelo autor (2026). Tokens de entrada incluem leitura de \textit{cache}. O tempo é o de parede do agente e não é comparável entre modelos.}
\end{table}

\textbf{Comportamento.} Houve um único \textit{timeout} válido (DeepSeek em
T05 \texttt{none}, que trabalhou os 30 minutos e passou mesmo assim) e nenhuma
contaminação. Nas 56 execuções de cada modelo, as regressões ficaram entre 1 e
5 para os modelos de API, contra 35 para o gpt-oss; nenhum agente editou
arquivos do \textit{harness} nem introduziu saída Blade sem escape; Sonnet e
Haiku alteraram testes preexistentes em 11 e 9 execuções, DeepSeek em 3 e Muse
em nenhuma. Três execuções tentaram buscar o repositório-sujeito no GitHub
(uma leitura de arquivo bruto, uma consulta à API e um \texttt{git ls-remote}),
todas do DeepSeek e todas no nível \texttt{none}; todas falharam com erro ou
404. Com $n = 3$ o fato é anedótico, mas é coerente com a ideia de que o
\textit{harness} dá ao agente contexto suficiente para não procurar fora.

%-------------------------------------------------------------------------------
\subsection{Resultados: qualidade de código}
\label{sec:resultados-qualidade}
\label{sec:qualidade-metodo}

A análise de qualidade cobriu as 309 execuções avaliadas da rodada (56 de cada
um dos cinco modelos completos e 29 do Beam), com as três camadas aplicadas a
todas; 306 execuções receberam as duas revisões de juízes e três ficaram com
uma, porque a segunda revisão terminou em erro do juiz. Os pares de juízes
foram fixos por modelo avaliado: Muse e DeepSeek para Sonnet e Haiku; Opus e
DeepSeek para o Muse; Opus e Muse para o DeepSeek; e rotação balanceada das
três duplas para gpt-oss e Beam. Cada juiz trabalhou em um espaço isolado com
o enunciado, o \textit{diff} sem os arquivos do \textit{harness} e sem os
\textit{lock files}, as seções normativas do \texttt{AGENTS.md} como padrão de
revisão (o mesmo para todos os níveis), a rubrica e o esquema de saída. Uma
verificação recusa qualquer espaço que contenha o identificador da execução, o
modelo, o nível ou o resultado dos testes. A revisão é validada contra um
esquema JSON e, se inválida, recebe uma segunda chance com os erros
informados.

\textbf{Efeito por nível.} A \autoref{tab:qualidade-efeito} resume o efeito de
cada nível contra o \texttt{none}, pareado por tarefa, nos 70 pares dos cinco
modelos completos. O \textit{harness} elevou a nota geral dos juízes em 0,31 a
0,34 ponto (em uma escala de 1 a 5; $p \le 0{,}001$), a dimensão de convenções
em 0,44 a 0,51 ($p < 0{,}001$) e a dimensão de testes em 0,56 a 0,82
($p < 0{,}001$). Na camada estática, o número de casos de teste novos subiu 10
a 11 por execução ($p < 0{,}001$), o escore de convenções subiu 0,07 a 0,10
($p \le 0{,}001$), os erros novos do Larastan caíram 0,5 a 0,6 por execução
($p$ entre 0,015 e 0,057) e a complexidade média caiu 0,16 a 0,19 ($p$ entre
0,001 e 0,051). Requisitos, segurança e escopo não mudaram ($p$ entre 0,06 e
0,82), nem a distância da solução de referência (Jaccard, $p \ge 0{,}17$). A
legibilidade subiu pouco (0,06 a 0,13) e a simplicidade não mudou de forma
consistente. O nível \texttt{skills} tem os maiores efeitos em convenções e
testes, mas as diferenças entre \texttt{agents}, \texttt{skills} e \texttt{full}
são pequenas diante dos efeitos de cada um contra o \texttt{none}: a maior parte
do ganho vem da primeira camada.

\begin{table}[htb]
    \centering
    \ABNTEXfontereduzida
    \caption{Efeito de cada nível contra \texttt{none} nas métricas de qualidade ($\Delta$ = média das diferenças pareadas por tarefa; Wilcoxon exato; 70 pares por nível, 63 a 70 nas métricas com dados ausentes)}
    \label{tab:qualidade-efeito}
    \begin{tabular}{lrrrrrr}
        \toprule
        & \multicolumn{2}{c}{\texttt{agents}} & \multicolumn{2}{c}{\texttt{skills}} & \multicolumn{2}{c}{\texttt{full}} \\
        \cmidrule(lr){2-3} \cmidrule(lr){4-5} \cmidrule(lr){6-7}
        \textbf{Métrica} & $\Delta$ & $p$ & $\Delta$ & $p$ & $\Delta$ & $p$ \\
        \midrule
        \multicolumn{7}{l}{\textit{Camada 1 e 2 (medidas objetivas)}} \\
        Escore de convenções (0 a 1) & $+0{,}07$ & 0,001 & $+0{,}10$ & $<0{,}001$ & $+0{,}07$ & $<0{,}001$ \\
        Larastan: erros novos & $-0{,}53$ & 0,034 & $-0{,}57$ & 0,057 & $-0{,}63$ & 0,015 \\
        Complexidade média & $-0{,}17$ & 0,001 & $-0{,}19$ & 0,025 & $-0{,}16$ & 0,051 \\
        Casos de teste novos & $+10{,}80$ & $<0{,}001$ & $+11{,}43$ & $<0{,}001$ & $+10{,}14$ & $<0{,}001$ \\
        Jaccard vs. referência & $-0{,}00$ & 0,265 & $-0{,}01$ & 0,245 & $-0{,}01$ & 0,169 \\
        \midrule
        \multicolumn{7}{l}{\textit{Camada 3 (juízes, notas de 1 a 5)}} \\
        Nota geral & $+0{,}31$ & 0,001 & $+0{,}34$ & 0,001 & $+0{,}33$ & $<0{,}001$ \\
        Requisitos & $+0{,}02$ & 0,313 & $-0{,}04$ & 0,732 & $+0{,}01$ & 0,815 \\
        Convenções & $+0{,}44$ & $<0{,}001$ & $+0{,}51$ & $<0{,}001$ & $+0{,}46$ & $<0{,}001$ \\
        Segurança & $+0{,}08$ & 0,230 & $+0{,}12$ & 0,056 & $+0{,}07$ & 0,171 \\
        Testes & $+0{,}67$ & $<0{,}001$ & $+0{,}82$ & $<0{,}001$ & $+0{,}56$ & $<0{,}001$ \\
        Simplicidade & $+0{,}13$ & 0,057 & $+0{,}04$ & 0,557 & $+0{,}09$ & 0,101 \\
        Legibilidade & $+0{,}13$ & 0,020 & $+0{,}06$ & 0,195 & $+0{,}11$ & 0,042 \\
        Escopo & $+0{,}01$ & 0,659 & $-0{,}01$ & 0,638 & $+0{,}03$ & 0,559 \\
        \bottomrule
    \end{tabular}
    \legend{Fonte: elaborado pelo autor (2026). O Beam, incompleto, não entra no agregado. $p$-valores sem correção para comparações múltiplas (39 comparações).}
\end{table}

\textbf{Efeito por modelo.} A \autoref{tab:qualidade-modelo} mostra a nota
geral dos juízes e o número de casos de teste novos por modelo e nível. O
efeito é maior em quem parte mais baixo. O DeepSeek v4.1 flash é quem mais
ganha: a nota geral vai de 3,18 no \texttt{none} para 3,93 a 4,11 com
\textit{harness}, a nota de testes de 2,79 para 4,43 a 4,57 e os casos de
teste de 8,0 para 27,8 a 32,1 por execução. O Muse Spark 1.3 vem em seguida:
nota geral de 2,61 para 3,04 a 3,32, testes de 1,21 para 2,07 a 2,75, casos de
teste de 1,1 para 5,4 a 9,7. Sonnet e Haiku ganham pouco na nota geral (de
3,96 para 4,14 a 4,17 e de 4,18 para 4,25 a 4,39), porque já partem de
convenções e testes altos, mas o Sonnet dobra os casos de teste escritos (de
16,7 para cerca de 34). Os dois modelos mais fracos não respondem ao
\textit{harness}: o gpt-oss-20b fica com nota geral entre 1,43 e 1,61 e nota de
testes de 1,0 em quase todos os níveis, e o Beam, nas células disponíveis, fica
entre 2,36 e 2,64, também sem testes. Em ambos, as regras do \texttt{AGENTS.md}
que mandam escrever testes simplesmente não são seguidas.

\begin{table}[htb]
    \centering
    \ABNTEXfontereduzida
    \caption{Nota geral dos juízes (média dos dois juízes, 1 a 5) e casos de teste novos por execução, por modelo e nível}
    \label{tab:qualidade-modelo}
    \begin{tabular}{lcccccccc}
        \toprule
        & \multicolumn{4}{c}{\textbf{Nota geral dos juízes}} & \multicolumn{4}{c}{\textbf{Casos de teste novos}} \\
        \cmidrule(lr){2-5} \cmidrule(lr){6-9}
        \textbf{Modelo} & \texttt{none} & \texttt{agents} & \texttt{skills} & \texttt{full} & \texttt{none} & \texttt{agents} & \texttt{skills} & \texttt{full} \\
        \midrule
        Sonnet 5.5 & 3,96 & 4,14 & 4,17 & 4,14 & 16,7 & 33,9 & 34,3 & 33,9 \\
        Haiku 5.5 & 4,18 & 4,39 & 4,31 & 4,25 & 36,8 & 45,3 & 45,4 & 42,0 \\
        Muse Spark 1.3 & 2,61 & 3,04 & 3,32 & 3,11 & 1,1 & 9,7 & 8,7 & 5,4 \\
        DeepSeek v4.1 flash & 3,18 & 4,11 & 3,93 & 4,07 & 8,0 & 27,8 & 28,4 & 32,1 \\
        Beam (free)$^{a}$ & 2,50 & 2,36 & 2,40 & 2,64 & 0,0 & 0,0 & 0,0 & 1,1 \\
        gpt-oss-20b & 1,61 & 1,43 & 1,46 & 1,61 & 0,0 & 0,0 & 2,9 & 0,0 \\
        \bottomrule
    \end{tabular}
    \legend{Fonte: elaborado pelo autor (2026). $n = 14$ por célula, exceto Sonnet \texttt{skills} ($n = 12$) e Haiku \texttt{skills} ($n = 13$) na nota geral. $^{a}$ Beam: $n$ = 10, 7, 5 e 7.}
\end{table}

\textbf{Um exemplo.} O caso do Muse Spark na tarefa T07 (chamada determinística
da próxima senha com modos de chamada e regra de justiça) ilustra o padrão. Nos
quatro níveis o modelo passou nos testes ocultos. No nível \texttt{none}, o
\textit{diff} trouxe a ação de domínio, o \textit{enum} de modos, a exceção e
a documentação, e nenhum arquivo de teste; os dois juízes deram nota 1 em
testes e nota geral 3, com a observação de que ``a lógica crítica de
concorrência é entregue sem nenhum teste''. No nível \texttt{skills}, o mesmo
modelo entregou, além do código, um arquivo de teste com quinze casos cobrindo
ordenação e empates, escopo por fila, os três modos e seus caminhos de erro,
fila vazia, a regra de justiça (disparo, reinício, histórico de outros modos) e
fila inativa; os juízes deram nota 4 em testes e nota geral 4. O oráculo oculto
não distingue as duas entregas; a revisão de código distingue.

\textbf{Concordância e validade dos juízes.} A \autoref{tab:kappa} mostra a
concordância entre as três duplas de juízes na nota geral. O $\kappa$ ponderado
vai de 0,71 (Muse e DeepSeek) a 0,85 (Opus e Muse), o que corresponde a
concordância substancial a quase perfeita na escala de
\citeonline{landis1977}; o $\kappa$ médio nas sete dimensões vai de 0,73 a 0,84.
Os juízes têm vieses pequenos e em direções distintas: o Muse é o mais
generoso (+0,22 em relação à média da dupla), o Opus o mais severo ($-0{,}17$) e
o DeepSeek fica perto do centro ($-0{,}09$). Como o par de juízes é fixo por
modelo avaliado, o viés do juiz é constante entre os níveis de um mesmo modelo
e não afeta as comparações pareadas, mas afeta a comparação de níveis absolutos
entre modelos.

\begin{table}[htb]
    \centering
    \ABNTEXfontereduzida
    \caption{Concordância entre juízes na nota geral e nas sete dimensões}
    \label{tab:kappa}
    \begin{tabular}{lcccc}
        \toprule
        \textbf{Dupla} & $n$ & $\kappa$ ponderado (geral) & Spearman (geral) & $\kappa$ médio (7 dimensões) \\
        \midrule
        Opus $\times$ Muse & 82 & 0,85 & 0,91 & 0,73 \\
        Opus $\times$ DeepSeek & 87 & 0,81 & 0,84 & 0,84 \\
        Muse $\times$ DeepSeek & 137 & 0,71 & 0,68 & 0,79 \\
        \bottomrule
    \end{tabular}
    \fonteautor
\end{table}

A \autoref{tab:validade} verifica se as notas medem algo externo a elas. A nota
geral correlaciona-se com o acerto nos testes ocultos ($\rho = 0{,}52$), e a
dimensão de requisitos um pouco mais ($\rho = 0{,}54$): os juízes, que não veem
o resultado dos testes, tendem a dar notas mais altas às mudanças que passam.
A dimensão de convenções dos juízes correlaciona-se fortemente com o escore
heurístico da camada 2 ($\rho = 0{,}73$), duas medidas independentes do mesmo
construto. Em contrapartida, a nota geral não se correlaciona com os erros novos
do Larastan ($\rho = -0{,}07$); no piloto, com duas tarefas, essa correlação era
de $-0{,}53$, e a divergência é discutida nas ameaças à validade.

\begin{table}[htb]
    \centering
    \ABNTEXfontereduzida
    \caption{Validade: correlação de Spearman entre medidas de camadas distintas (todas as execuções)}
    \label{tab:validade}
    \begin{tabular}{llcc}
        \toprule
        \textbf{Medida A} & \textbf{Medida B} & $\rho$ & $n$ \\
        \midrule
        Juízes: nota geral & Testes ocultos aprovados & 0,52 & 306 \\
        Juízes: requisitos & Testes ocultos aprovados & 0,54 & 306 \\
        Juízes: convenções & Escore de convenções (camada 2) & 0,73 & 305 \\
        Juízes: nota geral & Larastan: erros novos & $-0{,}07$ & 291 \\
        Escore de convenções (camada 2) & Testes ocultos aprovados & 0,37 & 308 \\
        \bottomrule
    \end{tabular}
    \fonteautor
\end{table}

\textbf{Regras de convenção.} Das dez regras heurísticas (Apêndice~C), as
mais violadas no conjunto foram C4 (escritas de domínio por classes de ação:
cumprida em 39 das 148 execuções em que se aplicou), C3 (autorização
\textit{server-side} de toda ação Livewire: 36 de 87), C2 (validação com
\texttt{\#[Validate]}: 53 de 106), C8 (\texttt{\$fillable} ou \texttt{\$guarded}
declarados: 84 de 151) e C1 (código de produção acompanhado de teste: 193 de
308). As regras de escape em Blade (C7), de ausência de camadas de serviço
(C6), de interface em Bulma (C9) e de migrações geradas e reversíveis (C5)
foram cumpridas em quase todas as execuções em que se aplicaram. No piloto,
as 25 violações apontadas foram lidas manualmente contra o código, e todas
eram reais (C1 19 de 19, C2 2 de 2, C10 4 de 4); uma correção na rodada
completa fez a regra C8 reconhecer a sintaxe de atributo \texttt{\#[Fillable]},
que antes era contada como violação.

%-------------------------------------------------------------------------------
\subsection{Discussão}
\label{sec:discussao}

\textbf{O \textit{harness} muda como o código é escrito, não se a tarefa é
resolvida.} Os dois conjuntos de resultados apontam na mesma direção. Pelo
oráculo oculto, nenhum nível alterou o \textit{pass@1} de nenhum modelo; pelos
juízes e pela camada estática, todos os níveis elevaram convenções, testes e
casos de teste escritos, sem tocar em requisitos, segurança ou escopo. A
dissociação é coerente com o que o \textit{harness} contém: o \texttt{AGENTS.md}
não ensina o domínio nem corrige a leitura de um enunciado, mas prescreve como
o código deve ser organizado e exige que cada mudança venha com testes. O
efeito sobre a dimensão de testes é, provavelmente, o mais direto: a regra
``nenhum código de produção sem teste'' é lida e seguida pelos modelos de
capacidade intermediária e alta. O exemplo da T07 mostra o limite do oráculo:
um PR sem nenhum teste e um PR com quinze testes passam igualmente nos testes
ocultos, mas um revisor humano não os aprovaria da mesma forma.

\textbf{Escala explica o \textit{pass@1}; o \textit{harness} explica parte da
qualidade.} Entre modelos, o \textit{pass@1} vai de 11\% a 100\%; entre níveis,
para um mesmo modelo, a variação é de no máximo 7 pontos percentuais. A escolha
do modelo explica quase toda a variação em acerto, e o \textit{harness} não
compensa a diferença: o Haiku 5.5 não alcança o Sonnet 5.5 em nenhum nível, e o
gpt-oss-20b não sai do patamar de 1 a 2 tarefas em 14. Na qualidade, o quadro é
diferente. O DeepSeek v4.1 flash com \textit{harness} recebe dos juízes uma nota
geral (3,93 a 4,11) comparável à do Sonnet 5.5 sem \textit{harness} (3,96) e
escreve um número de testes (28 a 32 por execução) semelhante ao do Sonnet
com \textit{harness} (34). Nesse sentido estrito, e apenas nessa métrica, a
estrutura de apoio aproximou um modelo intermediário de um modelo maior. A
leitura mais defensável do lema \textit{scaffolding over scale}, à luz destes
dados, é portanto condicional: a estrutura de apoio não substitui a escala na
capacidade de resolver a tarefa, mas pode substituir parte da escala na
disciplina com que o código é entregue, desde que o modelo seja capaz de ler
e seguir as regras.

\textbf{O efeito concentra-se no meio da distribuição.} Quem mais ganha são
os modelos que partem de uma posição intermediária (DeepSeek e Muse). Os
modelos de ponta já escrevem testes e seguem convenções sem instrução, e o
\textit{harness} acrescenta pouco. Os modelos mais fracos (gpt-oss-20b e Beam)
não respondem: as regras estão na árvore, mas o modelo não as lê ou não as
segue, e a nota de testes permanece em 1,0. Há, assim, um limiar de capacidade
abaixo do qual a estrutura de apoio é inerte, o que contraria a versão mais
forte da hipótese, segundo a qual o \textit{harness} compensaria justamente os
modelos menores.

\textbf{Qual camada carrega o efeito.} A maior parte do ganho de qualidade e
do acréscimo de custo aparece já no nível \texttt{agents}. O nível
\texttt{skills} acrescenta um pouco em convenções e testes (os maiores $\Delta$
estão nele), o que é coerente com o conteúdo das \textit{skills}, que detalham
exatamente essas práticas. O nível \texttt{full} não acrescenta nada mensurável
e, em alguns casos, devolve parte do ganho; a configuração de LSP, cuja intenção
é reduzir leituras, não reduziu o consumo de \textit{tokens} de entrada em
relação ao \texttt{skills} de forma consistente. Com uma repetição, as
diferenças entre os três níveis com \textit{harness} estão dentro do ruído, e a
conclusão segura é apenas que a primeira camada responde pela maior parte do
efeito.

\textbf{Custo contra benefício.} O \textit{harness} encarece cada execução em
dezenas de pontos percentuais em turnos e \textit{tokens} de entrada, para um
\textit{pass@1} idêntico. Se o critério for apenas resolver a tarefa medida por
testes, o \textit{harness} é custo sem retorno nesta amostra. Se o critério for
entregar um PR com testes e dentro das convenções do projeto, o acréscimo
compra um ganho de 0,3 ponto na nota geral e de 0,5 a 0,8 ponto em testes, mais
dez casos de teste por mudança, o que, em um fluxo real, reduziria o trabalho
de revisão e de complementação posterior. Parte do custo adicional é, aliás, o
custo de escrever e executar os testes que o \texttt{AGENTS.md} exige. A
relação custo-benefício depende, portanto, do que se pede ao agente e de quem
revisa o resultado.

\textbf{Relação com a literatura.} O resultado dialoga com o achado de
\citeonline{yang2024sweagent} de que a interface do agente altera a taxa de
resolução: aqui, a intervenção foi textual e versionada no repositório, e não na
ferramenta, e o seu efeito apareceu na qualidade, não na resolução, em um
conjunto de tarefas em que a resolução já estava no teto. A dissociação entre
correção funcional e qualidade ecoa \citeonline{liu2023evalplus} e
\citeonline{pearce2022asleep}, que mostram, por caminhos diferentes, que passar
nos testes não esgota o que se espera de um código. Por fim, a concordância
entre juízes e a sua correlação com medidas independentes estão em linha com
o que \citeonline{zheng2023judging} e \citeonline{zhuo2024icescore} relatam, e o
cuidado com a autopreferência segue \citeonline{panickssery2024self}.

%-------------------------------------------------------------------------------
\subsection{Ameaças à validade}
\label{sec:ameacas}

\textbf{Validade de conclusão.} Com uma repetição por célula e 14 tarefas por
modelo, o desenho detecta apenas efeitos grandes sobre o \textit{pass@1}: o
teste de McNemar exige ao menos seis discordâncias na mesma direção para
$p < 0{,}05$, de modo que efeitos de 5 a 15 pontos percentuais passariam
despercebidos. Ausência de evidência não é evidência de ausência. A variância
de um mesmo modelo na mesma célula não foi estimada; as 20 células de
repetições 2 e 3 da rodada completa cancelada ficam no manifesto como amostra
pequena dessa variação. Na análise de qualidade foram feitas 39 comparações sem
correção para múltiplos testes; com um controle de taxa de falsas descobertas
\cite{benjamini1995fdr}, os efeitos em convenções, testes, casos de teste e nota
geral ($p \le 0{,}001$) resistiriam, mas os de Larastan, complexidade e
legibilidade (0,015 a 0,057) devem ser lidos como sugestivos. As cinco
subtarefas dos dois PRs compostos são correlacionadas (T04b e T04c incluem os
testes de T04a; T03a e T03b compartilham rotas), o que reduz o número efetivo
de pares.

\textbf{Validade interna.} O efeito de teto é a principal ameaça: quatro
modelos resolvem 93\% ou mais das tarefas sem \textit{harness}, e os testes
ocultos pareceram pouco discriminativos para modelos fortes (fração aprovada de
0,993 ou mais mesmo nas falhas). Tarefas mais difíceis seriam necessárias para
medir ganho de \textit{pass@1} acima do teto. O Beam ficou incompleto (29 de 56
células, de dois \textit{endpoints} do mesmo modelo) e o Qwen foi retirado, de
modo que esses modelos não entram nas comparações. O limite de memória de 6~GB
foi introduzido durante a rodada; nenhuma execução posterior morreu por
memória, e a anterior afetada foi refeita. Três reinícios do orquestrador
recomeçaram agentes em andamento como novas tentativas, sem duplicar células.
O modelo local rodou quantizado e com o analisador de formato do llama.cpp;
24 das suas 56 falhas são de formato de saída, uma limitação do conjunto
modelo mais servidor e não apenas do modelo.

\textbf{Validade de construto.} O oráculo é parcialmente escrito com IA:
cerca de 15\% dos \textit{commits} do repositório-sujeito são rastreáveis a
ferramentas de IA e os testes ocultos foram, às vezes, escritos com
assistência. Os testes ocultos acoplam-se a nomes: uma implementação correta
com nomes diferentes falha, o que o contrato do enunciado mitiga sem eliminar.
Os juízes são LLMs, com os vieses conhecidos; a concordância entre eles é alta,
mas eles podem concordar no mesmo erro, e não houve calibração com revisores
humanos. O Muse Spark ajudou a escrever o código do repositório-sujeito, treina
com os dados enviados e é, ao mesmo tempo, modelo avaliado e juiz de outros
modelos; a checagem de modelo confirmou que apenas o modelo sob teste respondeu
em cada execução, e o Muse nunca julgou o próprio código, mas uma afinidade de
estilo com o repositório não pode ser descartada. A divergência entre Larastan
e juízes (correlação $-0{,}53$ no piloto e $-0{,}07$ na rodada completa) indica
que a contagem de erros novos do Larastan captura algo que os juízes não veem,
ou vice-versa; o Larastan não está instalado no projeto e os agentes não o
executam, e a aplicação não inicializa em parte das execuções do gpt-oss, o que
reduz o $n$ dessa métrica. As regras C1 a C10 são heurísticas: a validação
manual no piloto encontrou apenas violações reais, mas regras que não se
aplicaram no piloto (C5 e C8) só foram exercitadas na rodada completa, e uma
delas precisou de correção.

\textbf{Validade externa.} As tarefas vêm de um único repositório privado, de
um único autor, em uma única pilha (Laravel, Livewire, Bulma), e o
\textit{harness} é o desse repositório. Nada aqui autoriza generalizar para
outros repositórios, linguagens ou \textit{harnesses}. Os \textit{runners} são
diferentes entre famílias de modelos, de modo que comparações entre modelos
misturam modelo e \textit{runner}; apenas as comparações de níveis dentro de
cada modelo são controladas. O \textit{sandbox} tinha internet liberada; a
natureza privada do repositório e o detector cobrem o risco de vazamento, e
nenhum vazamento foi observado. O preâmbulo não interativo é igual para todos os
modelos e níveis, mas é uma condição específica deste experimento.
